% Standalone problem document, generated from the adjacent README.md.
% Compile with XeLaTeX. Mathematical content is maintained in Markdown.
\documentclass[11pt,a4paper]{article}
\usepackage[fontsize=10.5pt]{scrextend}
\usepackage[margin=22mm,headheight=15pt,headsep=9mm,footskip=11mm]{geometry}
\usepackage{fontspec}
\setmainfont{texgyrepagella}[Extension=.otf,UprightFont=*-regular,BoldFont=*-bold,ItalicFont=*-italic,BoldItalicFont=*-bolditalic]
\setsansfont{texgyreheros}[Extension=.otf,UprightFont=*-regular,BoldFont=*-bold,ItalicFont=*-italic,BoldItalicFont=*-bolditalic]
\setmonofont{texgyrecursor}[Extension=.otf,UprightFont=*-regular,BoldFont=*-bold,ItalicFont=*-italic,BoldItalicFont=*-bolditalic,Scale=MatchLowercase]
\usepackage{amsmath,amssymb}
\usepackage{unicode-math}
\setmathfont{texgyrepagella-math.otf}
\usepackage{xcolor,microtype,enumitem,needspace,fancyhdr}
\usepackage{bookmark}
\definecolor{ink}{HTML}{17344B}
\definecolor{accent}{HTML}{197D80}
\definecolor{muted}{HTML}{596773}
\hypersetup{colorlinks=true,linkcolor=accent,urlcolor=accent,pdftitle={AC-12 - Realizing
every sign-matrix permanent with negative entries confined on or above
the diagonal},pdfauthor={Open Problems in Numerical Linear Algebra}}
\urlstyle{same}
\setlength{\parindent}{0pt}
\setlength{\parskip}{5pt plus 1pt minus 1pt}
\linespread{1.04}
\setlength{\emergencystretch}{3em}
\widowpenalty=10000
\clubpenalty=10000
\displaywidowpenalty=10000
\allowdisplaybreaks[1]
\setlist{nosep,leftmargin=1.5em}
\providecommand{\tightlist}{\setlength{\itemsep}{0pt}\setlength{\parskip}{0pt}}
\providecommand{\pandocbounded}[1]{#1}
\pagestyle{fancy}
\fancyhf{}
\fancyhead[L]{\sffamily\footnotesize\color{muted}OPEN PROBLEMS IN NUMERICAL LINEAR ALGEBRA}
\fancyhead[R]{\sffamily\bfseries\footnotesize\color{accent}AC-12}
\fancyfoot[L]{\parbox[b]{0.9\textwidth}{\sffamily\fontsize{8}{10}\selectfont\color{muted}Editorial ratings; a bounded search cannot certify that a problem remains unsolved.\\Literature check: 2026-09-10}}
\fancyfoot[R]{\sffamily\footnotesize\color{muted}\thepage}
\renewcommand{\headrulewidth}{0.3pt}
\renewcommand{\headrule}{\hbox to\headwidth{\color{accent}\leaders\hrule height \headrulewidth\hfill}}
\makeatletter
\renewcommand{\section}{\Needspace{5\baselineskip}\@startsection{section}{1}{0pt}{12pt plus 2pt minus 2pt}{4pt}{\sffamily\bfseries\large\color{ink}}}
\renewcommand{\subsection}{\Needspace{4\baselineskip}\@startsection{subsection}{2}{0pt}{9pt plus 1pt minus 1pt}{3pt}{\sffamily\bfseries\normalsize\color{ink}}}
\makeatother
\setcounter{secnumdepth}{0}
\begin{document}
{\sffamily\small\color{muted}Arithmetic and complexity\par}
\vspace{3pt}
{\raggedright\hyphenpenalty=10000\exhyphenpenalty=10000\sffamily\bfseries\fontsize{21}{24}\selectfont\color{ink}Realizing
every sign-matrix permanent with negative entries confined on or above
the diagonal\par}
\vspace{7pt}
{\sffamily\small\textbf{Difficulty:} challenging\quad\textbf{Importance:} interesting
to specialist\par}
{\sffamily\small\bfseries\color{accent}Status: Open\par}
\vspace{4pt}
{\color{accent}\hrule height 0.8pt}
\vspace{6pt}
\textbf{Rating rationale:} Challenging because equality of the full
permanent ranges requires a uniform value-preserving construction or an
obstruction; specialist importance concerns structural realization of
discrete matrix invariants.

\subsection{Problem statement}\label{problem-statement}

For each integer \(n\ge1\), let \[
\Omega_n=\{-1,1\}^{n\times n},\qquad
\mathcal U_n=\{B\in\Omega_n:b_{ij}=1\text{ whenever }i>j\}.
\] Define
\(\operatorname{per}A=\sum_{\sigma\in S_n}\prod_{i=1}^n a_{i,\sigma(i)}\).
Does every \(A\in\Omega_n\) admit some \(B\in\mathcal U_n\) with \[
\operatorname{per}B=\operatorname{per}A?
\] Equivalently, do \(\Omega_n\) and \(\mathcal U_n\) have identical
permanent ranges for every \(n\)? Entries below the diagonal are \(+1\),
not zero. Negative entries are allowed on the diagonal.

\subsection{Relevance and ratings}\label{relevance-and-ratings}

A positive answer would reduce exact value realization to a smaller,
explicitly structured family of matrices. The question connects matrix
structure and the combinatorial complexity of the permanent; it does not
ask for an efficient permanent-evaluation algorithm. The main
significance is to specialists, and a uniform structural argument or
counterexample is required.

\subsection{References}\label{references}

\begin{itemize}
\tightlist
\item
  D. Ingram and A. Razborov, \emph{On the range of the permanent of
  \((\pm1)\)-matrices}, Linear Algebra Appl. 743 (2026), 271--285
  (\href{https://doi.org/10.1016/j.laa.2026.04.027}{journal});
  \href{https://arxiv.org/html/2507.09433v1}{arXiv:2507.09433v1}, §6,
  Problem 4 and the paragraph defining the restricted family.
\item
  Z. Hunter, M. Kwan and L. Sauermann, \emph{Exponential
  anticoncentration of the permanent} (2025), §1.2, Corollary 1.3
  (\href{https://arxiv.org/abs/2509.22577}{primary paper}).
\end{itemize}

\subsection{Status check --- 2026-09-10}\label{status-check-2026-09-10}

Checked \href{https://arxiv.org/html/2507.09433v1}{Ingram--Razborov, §6,
Problem 4}, including the preceding definition: negative entries may
occur on or above the diagonal. The displayed target matches that
question. Exact-title and upper-triangular permanent-range searches
found no general resolution.
\href{https://arxiv.org/abs/2509.22577}{Hunter--Kwan--Sauermann,
Corollary 1.3}, settles exponential growth of the unrestricted range,
not equality with this restricted family; rank-dependent maximum results
also leave exact ranges undetermined.
\end{document}
